% How one market model becomes 373 names (user, 2026-10-30: mock-up B, a flow), for the slide after the JumpHMM
% slide. Input once in the preamble, after vnslides.
% Each name's daily growth rate is the sum of four parts (design: docs/superpowers/specs/
% 2026-10-26-synthetic-market-stress-test-design.md §4; code: compose_name in code/src/SyntheticMarket.jl):
%   - the trend α_i, the name's 2014–2024 SIM α pulled toward the universe mean (φ = 0.6 in
%     results/synthetic/generator_summary.toml);
%   - β_i(t) g_m(t): the market JumpHMM, fit to SPY 2014–2024, the generator on the JumpHMM slide; β_i drifts as an
%     EMA of random shocks with a half-life of 126 trading days (h_beta), six months;
%   - δ_i f_s(t): one JumpHMM per sector, fit to the sector's mean SIM residual, 11 in all; δ_i is the name's loading;
%   - s_i ε̃_i(t): one JumpHMM per name, fit to its own 2014–2024 growth, 373 in all; s_i holds the name's variance at
%     its own model's level.
% The paths run over the real calendar, 2014-01-03 to 2026-09-04.
%
% \vnmultiasset draws the flow: the four inputs as boxes in a column (the trend dashed, as it is a fixed number, not a
% model; the SPY model gold, as SPY is everywhere; the sector models link blue), an arrow from each into one sum bar
% with its term above the arrow and a short note below, and the sum's output, g_i(t). Every box is one width, so the
% arrows start level.
\newcommand{\vnmultiasset}{%
\begin{tikzpicture}[
    gen/.style={draw, line width=1pt, rounded corners=2pt, text width=4.6cm, align=flush left, inner xsep=6pt,
      inner ysep=2.5pt, anchor=west, fill=vnbg, font=\fontsize{7.5}{9}\selectfont, text=vnslidebody},
    flow/.style={->, >={Stealth[length=5pt, width=4.5pt]}, draw=vnslidebody, line width=0.8pt},
    term/.style={font=\small, text=vnslideink, inner sep=0pt},
    note/.style={font=\fontsize{7.5}{9}\selectfont, text=vnslidebody, inner sep=0pt},
  ]
  % the inputs, one row each, 1.15 cm apart
  \node[gen, draw=vnslidebody, dash pattern=on 3pt off 2pt] (ga) at (0, 1.725)
    {\vnmultiassetname{Trend}\\[1.5pt]2014–2024 value, shrunk to the mean};
  \node[gen, draw=vnfigaccent] (gm) at (0, 0.575)
    {\vnmultiassetname{SPY JumpHMM}\\[1.5pt]1 model, fit to SPY 2014–2024};
  \node[gen, draw=vnslidelink] (gs) at (0, -0.575)
    {\vnmultiassetname{Sector JumpHMMs}\\[1.5pt]11 models, fit to sector residuals};
  \node[gen, draw=vnslideink!70!vnbg] (ge) at (0, -1.725)
    {\vnmultiassetname{Name JumpHMMs}\\[1.5pt]373 models, fit to each name’s growth};
  % the sum, a bar the four arrows run into
  \node[draw=vnslidebody, line width=0.8pt, rounded corners=2pt, minimum width=0.5cm, minimum height=4.3cm,
    inner sep=0pt, fill=vnbg, text=vnslideink] (sum) at (8.9, 0) {$+$};
  % each arrow carries its term above it and a note below it
  \draw[flow] (ga.east) -- (ga.east -| sum.west) coordinate[midway] (vnmaA);
  \node[term, anchor=south] at ([yshift=3pt]vnmaA) {$\alpha_i$};
  \draw[flow] (gm.east) -- (gm.east -| sum.west) coordinate[midway] (vnmaM);
  \node[term, anchor=south] at ([yshift=3pt]vnmaM) {$\beta_i(t)\, g_m(t)$};
  \node[note, anchor=north] at ([yshift=-3.5pt]vnmaM) {$\beta_i$ drifts, half-life six months};
  \draw[flow] (gs.east) -- (gs.east -| sum.west) coordinate[midway] (vnmaS);
  \node[term, anchor=south] at ([yshift=3pt]vnmaS) {$\delta_i\, f_s(t)$};
  \node[note, anchor=north] at ([yshift=-3.5pt]vnmaS) {$\delta_i$: sector loading};
  \draw[flow] (ge.east) -- (ge.east -| sum.west) coordinate[midway] (vnmaE);
  \node[term, anchor=south] at ([yshift=3pt]vnmaE) {$s_i\, \tilde\varepsilon_i(t)$};
  \node[note, anchor=north] at ([yshift=-3.5pt]vnmaE) {$s_i$ holds its variance};
  % the output
  \draw[flow] (sum.east) -- ++(0.7, 0) node[term, anchor=west, xshift=3pt] (gi) {$g_i(t)$};
  \node[note, anchor=north west, align=left] at ([yshift=-4pt]gi.south west)
    {for each of\\373 names,\\2014 – Sep 2026};
\end{tikzpicture}%
}

% a box's first line: the input's name, a size up and in the heading color
\newcommand{\vnmultiassetname}[1]{{\fontsize{9.5}{11}\selectfont\vnmedium\color{vnslideink}#1}}
